% Part: first-order-logic
% Chapter: models-theories
% Section: theories

\documentclass[../../../include/open-logic-section]{subfiles}

\begin{document}

\olfileid{fol}{mat}{the}

\olsection{प्रथम-क्रम उपपत्तींची उदाहरणे}

\begin{ex}
भाषा~$\Lang L_<$ मधील काटेकोर रेषीय क्रमांची उपपत्ती पुढील संचाने
स्वयंसिद्धकांद्वारे निरूपित होते:
\begin{align*}
\{\quad & \lforall[x][\lnot x < x], \\
& \lforall[x][\lforall[y][((x < y \lor y <
    x) \lor x = y)]], \\
& \lforall[x][\lforall[y][\lforall[z][((x < y
      \land y < z) \lif x < z)]]] \quad \}
\end{align*}
ही उपपत्ती अभिप्रेत !!{structure}s पूर्णपणे पकडते: प्रत्येक काटेकोर
रेषीय क्रम या स्वयंसिद्धक-प्रणालीचे प्रतिमान असतो; आणि उलट, $R$ हा
संच~$X$ वरील रेषीय क्रम असेल, तर $\Domain M = X$ आणि
$\Assign{<}{M} = R$ असलेली रचना~$\Struct M$ या उपपत्तीचे प्रतिमान असते.
\end{ex}

\begin{ex}
$\Obj 1$ (!!{constant}) आणि $\cdot$ (द्विस्थानी !!{function}) असलेल्या
भाषेतील गटांची उपपत्ती पुढील वाक्यांनी स्वयंसिद्धकांद्वारे निरूपित होते:
\begin{align*}
& \lforall[x][\eq[(x \cdot \Obj 1)][x]]\\
& \lforall[x][\lforall[y][\lforall[z][\eq[(x \cdot (y \cdot z))][((x
          \cdot y) \cdot z)]]]]\\
& \lforall[x][\lexists[y][\eq[(x \cdot y)][\Obj 1]]]
\end{align*}
\end{ex}

\begin{ex}
पेआनो अंकगणिताची उपपत्ती अंकगणिताची भाषा~$\Lang L_A$ मधील पुढील
वाक्यांनी स्वयंसिद्धकांद्वारे निरूपित होते.
\begin{align*}
& \lforall[x][\lforall[y][(\eq[x'][y'] \lif \eq[x][y])]]\\
& \lforall[x][\eq/[\Obj 0][x']]\\
& \lforall[x][\eq[(x + \Obj 0)][x]]\\
& \lforall[x][\lforall[y][\eq[(x + y')][(x + y)']]]\\
& \lforall[x][\eq[(x \times \Obj 0)][\Obj 0]]\\
& \lforall[x][\lforall[y][\eq[(x \times y')][((x \times y) + x)]]]\\
& \lforall[x][\lforall[y][(x < y \liff \lexists[z][\eq[(z' + x)][y])]]]\\
\intertext{आणि पुढील रूपाची सर्व वाक्ये}
& (!A(\Obj 0) \land \lforall[x][(!A(x) \lif !A(x'))]) \lif \lforall[x][!A(x)]
\end{align*}
शेवटच्या रूपाची अनंतपणे अनेक वाक्ये असल्यामुळे ही स्वयंसिद्धक-प्रणाली
अनंत आहे. शेवटच्या रूपाला \emph{विगमन-रूपबंध} म्हणतात. (प्रत्यक्षात,
विगमन-रूपबंधाचे स्वरूप आपण येथे दाखवले आहे त्यापेक्षा थोडे अधिक गुंतागुंतीचे
आहे.)

शेवटचे स्वयंसिद्धक हे~$<$ ची \emph{स्पष्ट व्याख्या} आहे.
\end{ex}

\begin{ex}
शुद्ध संचांची उपपत्ती गणिताच्या पायामध्ये (आणि गणिताच्या तत्त्वज्ञानातही)
महत्त्वाची भूमिका बजावते. एखाद्या संचाचे सर्व !!{element}s हेही शुद्ध संच
असतील, तर तो संच शुद्ध असतो. म्हणून रिक्त संच शुद्ध मानला जातो; पण एखाद्या
संचात संच नसलेली कोणतीही वस्तू !!a{element} म्हणून असेल, तर तो संच शुद्ध
नसतो. म्हणजे शुद्ध संच हे फक्त रिक्त संचापासून, कोणतेही ``मूलघटक'' न
वापरता बनवलेले संच होत; मूलघटक म्हणजे स्वतः संच नसलेल्या वस्तू.

पुढील वाक्ये शुद्ध संचांच्या उपपत्तीसाठी स्वयंसिद्धक-प्रणाली मानता येतील:
\begin{align*}
& \lexists[x][\lnot \lexists[y][y \in x]]\\
& \lforall[x][\lforall[y][(\lforall[z](z \in x \liff z \in y) \lif
\eq[x][y])]]\\
& \lforall[x][\lforall[y][\lexists[z][\lforall[u][(u \in z \liff
(\eq[u][x] \lor \eq[u][y]))]]]]\\
& \lforall[x][\lexists[y][\lforall[z][(z \in y \liff \lexists[u][(z \in
        u \land u \in x)])]]]\\
\intertext{आणि पुढील रूपाची सर्व वाक्ये} &
\lexists[x][\lforall[y][(y \in x \liff !A(y))]]
\end{align*}
पहिले स्वयंसिद्धक सांगते की कोणतेही !!{element}s नसलेला एक संच आहे (म्हणजे
$\emptyset$ अस्तित्वात आहे); दुसरे सांगते की संच विस्तारात्मक आहेत; तिसरे
सांगते की कोणत्याही संच~$X$ आणि~$Y$ साठी संच~$\{X, Y\}$ अस्तित्वात आहे;
आणि चौथे सांगते की कोणत्याही संच~$X$ साठी संच~$\cup X$ अस्तित्वात आहे,
जिथे~$\cup X$ हा~$X$ च्या सर्व घटकांचा संयोग आहे.

सर्वांत शेवटी उल्लेखलेल्या !!{sentence}s ना एकत्रितपणे
\emph{अनौपचारिक गुणधर्माधारित संचरचना-रूपबंध} म्हणतात. तो मूलतः असे सांगतो
की प्रत्येक~$!A(x)$ साठी संच~$\Setabs{x}{!A(x)}$ अस्तित्वात आहे---म्हणून
पहिल्या नजरेत ते सत्य, उपयुक्त आणि कदाचित आवश्यकही वाटणारे स्वयंसिद्धक आहे.
त्याला ``अनौपचारिक'' म्हणतात कारण, प्रत्यक्षात, त्यामुळे ही उपपत्ती
अपूर्ततायोग्य होते: $!A(y)$ म्हणून $\lnot y \in y$ घेतल्यास पुढील
!!{sentence} मिळते:
\[
\lexists[x][\lforall[y][(y \in x \liff \lnot y \in y)]]
\]
आणि या !!{sentence} ची कोणत्याही !!{structure} मध्ये पूर्ति होत नाही.
\end{ex}

\begin{ex}
\emph{अवयवमीमांसेमध्ये} \emph{अवयवत्वाचा} संबंध हा एक मूलभूत संबंध आहे.
संचांच्या उपपत्तींप्रमाणेच अवयवत्वाच्याही अशा उपपत्ती आहेत, ज्या या
संबंधाविषयीच्या विविध (आणि कधीकधी परस्परविरोधी) संकल्पना
स्वयंसिद्धकांद्वारे निरूपित करतात.

अवयवमीमांसेच्या भाषेत एकच द्विस्थानी विधेयचिन्ह~$\Obj P$ असते आणि
$\Atom{\Obj P}{x, y}$ चा ``अर्थ'' $x$ हा~$y$ चा अवयव आहे असा असतो. हे
अर्थनिर्धारण मनात ठेवले असता, या भाषेतील !!a{structure} ला
\emph{अवयवत्वाची रचना} म्हणतात. अर्थात, एकच द्विस्थानी विधेय असलेली प्रत्येक
रचना या नावास खरोखर पात्र असेलच असे नाही. ``अवयवत्व'' पकडण्याची शक्यता
असण्यासाठी $\Assign{\Obj P}{M}$ ने काही अटींची पूर्ति केली पाहिजे; त्या अटी
आपण अवयवत्वाच्या उपपत्तीची स्वयंसिद्धके म्हणून मांडू शकतो. उदाहरणार्थ,
वस्तूंवरील अवयवत्व हा अंशतः क्रम आहे: प्रत्येक वस्तू स्वतःचाच अवयव असते
(जरी तो \emph{अनुचित} अवयव असला तरी); कोणत्याही दोन भिन्न वस्तू एकमेकींचे
अवयव असू शकत नाहीत; आणि एखाद्या वस्तूच्या अवयवाचा अवयव हा स्वतः त्या
वस्तूचा अवयव असतो. लक्षात घ्या की या अर्थाने ``चा अवयव आहे'' हे ``चा
उपसंच आहे'' यासारखे आहे; पण ते ``चा घटक आहे'' यासारखे नाही, कारण घटक-संबंध
परावर्तीही नाही आणि संक्रमकही नाही.
\begin{align*}
& \lforall[x][\Atom{\Obj P}{x,x}] \\
& \lforall[x][\lforall[y][((\Part{x}{y} \land \Part{y}{x})
      \lif \eq[x][y])]] \\
& \lforall[x][\lforall[y][\lforall[z][((\Part{x}{y} \land
        \Part{y}{z}) \lif \Part{x}{z})]]]\\
\intertext{याखेरीज, कोणत्याही दोन वस्तूंचा अवयवी संयोग असतो (ज्याचे त्या दोन
  वस्तू अवयव असतात आणि जो या बाबतीत न्यूनतम असतो).}  &
\lforall[x][\lforall[y][\lexists[z][\lforall[u][(\Part{z}{u} \liff
        (\Part{x}{u} \land \Part{y}{u}))]]]]
\end{align*}
तत्त्वमीमांसकांनी विचारात घेतलेल्या अवयवत्वाच्या मूलभूत तत्त्वांपैकी ही
फक्त काही तत्त्वे आहेत. मात्र, आधी काही परिभाषित संबंध न मांडता पुढील
तत्त्वांची सूत्रे बनवणे किंवा ती लिहून काढणे लवकरच कठीण होते. उदाहरणार्थ,
अवयवमीमांसेत रस असलेले बहुतेक तत्त्वमीमांसक पुढील तत्त्वही वैध मानतात:
जेव्हा वस्तू~$x$ ला उचित अवयव~$y$ असतो, तेव्हा तिला असा अवयव~$z$ देखील
असतो की त्या अवयवाचा~$y$ शी कोणताही सामाईक अवयव नसतो आणि त्यामुळे~$y$ व
$z$ यांचा अवयवी संयोग~$x$ असतो.
\end{ex}

\end{document}
